% Irodov Part VI §6.2 original problems 49--65.
% Source: supplied Arihant PDF pp.247--248, printed pp.238--239.
% Answers PDF pp.349--350. Every TeX derivation has an excheck clause.
\wavequantproblem{49}{238}
动能均为 $T=100\,\mathrm{eV}$ 的电子、
质子和铀原子，分别求其德布罗意波长。
\exsolution
三个粒子均在非相对论速度区间，
$p=\sqrt{2mT}$，所以
\[
 \boxed{\lambda_{\rm dB}=\frac h{\sqrt{2mT}}}.
\]
代
$m_e=9.109\times10^{-31}\,\mathrm{kg}$、
$m_p=1.673\times10^{-27}\,\mathrm{kg}$、
$m_{\rm U}\simeq238u$，得到
\[
 \boxed{\lambda_e\simeq123\,\mathrm{pm}},\qquad
 \boxed{\lambda_p\simeq2.86\,\mathrm{pm}},\qquad
 \boxed{\lambda_{\rm U}\simeq0.186\,\mathrm{pm}}.
\]
\excheck 相同动能下 $\lambda\propto m^{-1/2}$；
最重的铀原子波长最短。

\wavequantproblem{50}{238}
一个电子的德布罗意波长从
$100\,\mathrm{pm}$ 降到 $50\,\mathrm{pm}$，
问需要增加多少动能？
\exsolution
在这一能区电子可采用非相对论关系
\[
 T(\lambda)=\frac{h^2}{2m_e\lambda^2}.
\]
$\lambda_2=\lambda_1/2$，所以
$T_2=4T_1$，新增动能为
\[
 \boxed{\Delta T=T_2-T_1=
 \frac{3h^2}{2m_e(100\,\mathrm{pm})^2}
 \simeq0.45\,\mathrm{keV}}.
\]
\excheck 波长减半使动量加倍、动能增至4倍，
但\emph{增加}的是初动能的3倍。

\wavequantproblem{51}{238}
动能为 $T=25\,\mathrm{eV}$ 的中子撞击
起初静止的氘核。
求两粒子在质心参考系内各自的德布罗意波长。
\exsolution
实验室系中中子速度
$v_n=\sqrt{2T/m_n}$。
质心速度
$V_{\rm CM}=m_nv_n/(m_n+m_d)$。
中子质心动量模
\[
 p_n^*=m_n(v_n-V_{\rm CM})
 =\frac{m_nm_d}{m_n+m_d}v_n=\mu v_n.
\]
氘核的质心动量与它等大反向，
所以两者波长相同：
\[
 \boxed{\lambda_n^*=\lambda_d^*=
 \frac{h(m_n+m_d)}
 {m_d\sqrt{2m_nT}}\simeq8.6\,\mathrm{pm}}.
\]
\excheck 质心系总动量为0，
因此两粒子动量模相等，尽管质量不等。

\wavequantproblem{52}{238}
两相同非相对论粒子以相互垂直的方向运动，
在实验室系的德布罗意波长为
$\lambda_1,\lambda_2$。
求它们在质心系内各自的德布罗意波长。
\exsolution
质量均为 $m$，实验室动量
$\boldsymbol p_1\perp\boldsymbol p_2$，
$p_i=h/\lambda_i$。
质心速度
$\boldsymbol V=(\boldsymbol p_1+\boldsymbol p_2)/(2m)$。
第一粒子质心动量
$\boldsymbol p'_1=\boldsymbol p_1-m\boldsymbol V
=(\boldsymbol p_1-\boldsymbol p_2)/2$，
第二粒子为它的相反数，大小
$|\boldsymbol p'_1|=\tfrac12\sqrt{p_1^2+p_2^2}$。
于是
\[
 \boxed{\lambda^*=
 \frac{2}{\sqrt{1/\lambda_1^2+1/\lambda_2^2}}
 =\frac{2\lambda_1\lambda_2}
 {\sqrt{\lambda_1^2+\lambda_2^2}}}.
\]
\excheck 只有因为两实验室动量\emph{垂直}，
平方和才没有交叉项。

\wavequantproblem{53}{238}
室温下氢分子的\emph{最概然速率}对应
的德布罗意波长是多少？
\exsolution
麦克斯韦速度分布
$f(v)\propto v^2\exp[-mv^2/(2k_{\mathrm B}T)]$，
其速率峰值
$v_{\rm mp}=\sqrt{2k_{\mathrm B}T/m}$。
代入 $\lambda=h/(mv)$，
\[
 \boxed{\lambda(v_{\rm mp})=
 \frac{h}{\sqrt{2mk_{\mathrm B}T}}}.
\]
氢分子质量约 $m=2.016u$，
取 $T=300\,\mathrm K$，
\[
 \boxed{\lambda\simeq128\,\mathrm{pm}}.
\]
\excheck 本题问的是\emph{速率分布峰值对应的波长}，
与下一题“波长分布本身的峰值”不是同一概念。

\wavequantproblem{54}{238}
处于室温热平衡的氢分子，
若直接统计它们的德布罗意\emph{波长分布}，
求最概然波长。
\exsolution
设 $\lambda=h/(mv)$。
麦克斯韦速率分布
$f(v)\,dv\propto
v^2e^{-mv^2/(2k_{\mathrm B}T)}\,dv$。
进行概率变量替换：
\[
 \varphi(\lambda)\,d\lambda
 =f\!\left(\frac h{m\lambda}\right)
 \left|\frac{dv}{d\lambda}\right|d\lambda
 \propto\lambda^{-4}
 \exp\!\left[-\frac{h^2}
 {2mk_{\mathrm B}T\lambda^2}\right]d\lambda.
\]
极值条件 $d(\ln\varphi)/d\lambda=0$：
$-4/\lambda+h^2/(mk_{\mathrm B}T\lambda^3)=0$，
从而
\[
 \boxed{\lambda_{\rm mp}=
 \frac h{2\sqrt{mk_{\mathrm B}T}}
 \simeq0.089\,\mathrm{nm}\approx0.09\,\mathrm{nm}}.
\]
\excheck 概率密度在非线性变量变换下会改变形状；
不能把速率峰值代入 $\lambda=h/(mv)$ 当成波长峰值。

\wavequantproblem{55}{238}
推导动能为 $T$ 的相对论粒子的德布罗意波长；
若使用非相对论公式，要求相对误差不超过
$1\%$，分别求电子和质子的动能上限。
\exsolution
总能量 $E=mc^2+T$，
由相对论能量动量关系
$E^2=p^2c^2+m^2c^4$，得
\[
 pc=\sqrt{T(T+2mc^2)},\quad
 \boxed{\lambda_{\rm rel}=
 \frac{hc}{\sqrt{T(T+2mc^2)}}}.
\]
非相对论近似
$\lambda_{\rm nr}=h/\sqrt{2mT}$，
因此
\[
 \frac{\lambda_{\rm nr}}{\lambda_{\rm rel}}
 =\sqrt{1+\frac T{2mc^2}}.
\]
限定
$\lambda_{\rm nr}/\lambda_{\rm rel}-1\le0.01$，
有
\[
 \boxed{T\le2mc^2[(1.01)^2-1]
 =0.0402\,mc^2}.
\]
故
\[
 \boxed{T_{e,\max}\simeq20.5\,\mathrm{keV}},\qquad
 \boxed{T_{p,\max}\simeq37.7\,\mathrm{MeV}}.
\]
\excheck 1%是对\emph{波长}的相对误差；
非相对论波长比精确波长稍大。

\wavequantproblem{56}{238}
电子的德布罗意波长等于其普通
康普顿波长 $\lambda_C=h/(m_ec)$ 时，
其动能应为多少？
\exsolution
德布罗意式 $\lambda=h/p$，
令其等于 $h/(m_ec)$，得
$p=m_ec$。
由
$E=\sqrt{m_e^2c^4+p^2c^2}=\sqrt2m_ec^2$，
故动能
\[
 \boxed{T=(\sqrt2-1)m_ec^2
 \simeq0.212\,\mathrm{MeV}}.
\]
\excheck 所谓康普顿波长是 $h/(mc)$，
不是约化形式 $\hbar/(mc)$。

\wavequantproblem{57}{238}
X 射线管连续谱最短波长为
$\lambda_{\rm sh}=10.0\,\mathrm{pm}$。
求电子撞击阳极前的相对论德布罗意波长。
\exsolution
能量守恒给电子动能
$T=hc/\lambda_{\rm sh}$。
用相对论动量
$pc=\sqrt{T(T+2m_ec^2)}$，
于是
\[
 \boxed{\lambda_e=
 \frac{hc}{\sqrt{T(T+2m_ec^2)}}
 =\frac{\lambda_{\rm sh}}
 {\sqrt{1+\dfrac{2m_ec^2\lambda_{\rm sh}}{hc}}}}.
\]
数值 $T\simeq124\,\mathrm{keV}$，
\[
 \boxed{\lambda_e\simeq3.3\,\mathrm{pm}}.
\]
\excheck X 射线短波限是光子的波长，
不等于撞靶\emph{电子}的德布罗意波长。

\wavequantproblem{58}{238}
动能单一的电子束正入射宽
$b=1.0\,\mu\mathrm m$ 的窄方缝，
距缝 $\ell=50\,\mathrm{cm}$ 的屏上
单缝衍射中央主极大宽
$\Delta x=0.36\,\mathrm{mm}$。
求电子速度。
\exsolution
单缝第一暗纹满足
$b\sin\theta_1=\lambda_{\rm dB}$。
小角下 $\theta_1\simeq\lambda_{\rm dB}/b$，
中央主极大从左右第一暗纹之间的总宽度为
$\Delta x\simeq2\ell\lambda_{\rm dB}/b$，
即
\[
 \lambda_{\rm dB}=\frac{b\Delta x}{2\ell}.
\]
因 $\lambda_{\rm dB}=h/(m_ev)$，
\[
 \boxed{v=\frac{2h\ell}{m_eb\Delta x}
 \simeq2.0\times10^6\,\mathrm{m/s}}.
\]
\excheck 中央主极大宽度是从 $-\theta_1$ 到
$+\theta_1$ 的\emph{全宽}，不能少因子2。

\wavequantproblem{59}{238}
电子束经电压 $U=25\,\mathrm V$ 加速，
正入射两细缝，间距
$d=50\,\mu\mathrm m$。
屏距 $\ell=100\,\mathrm{cm}$，
求干涉条纹相邻亮纹距离。
\exsolution
加速电子获得动能
$T=eU$，非相对论动量
$p=\sqrt{2m_eeU}$，故波长
$\lambda=h/\sqrt{2m_eeU}$。
双缝相邻干涉亮纹角距满足
$\Delta\theta\simeq\lambda/d$，
屏上线距
\[
 \boxed{\Delta x=
 \frac{\ell h}{d\sqrt{2m_eeU}}
 \simeq4.9\,\mu\mathrm m}.
\]
\excheck 双缝条纹间距与双缝间距成反比，
且随加速电压按 $U^{-1/2}$ 缩小。

\wavequantproblem{60}{239}
动能单一的电子束以
$\theta=30^\circ$ 掠入射铝单晶天然晶面。
相邻晶面距 $d=0.20\,\mathrm{nm}$，
反射极大在某加速电压 $U_0$ 时出现。
若电压增为原来的
$\eta=2.25$ 倍，又出现相邻下一级反射极大，
求 $U_0$。
\exsolution
布拉格反射条件 $2d\sin\theta=m\lambda_m$。
非相对论电子
$\lambda\propto U^{-1/2}$，
因此电压变为 $\eta U_0$ 时
$\lambda_{\rm new}=\lambda_0/\sqrt\eta$。
对同一晶面和入射角，
$\lambda\propto1/m$，相邻两个整数级次满足
\[
 \frac{m+1}{m}=\sqrt\eta=1.5,
 \quad\Longrightarrow\quad m=2.
\]
所以
$2d\sin\theta=2\lambda_0$，
即 $\lambda_0=d\sin\theta$。
于是
\[
 \boxed{U_0=\frac{h^2}
 {2m_e e\,d^2\sin^2\theta}
 \simeq0.15\,\mathrm{kV}=150\,\mathrm V}.
\]
\excheck 电压提高使波长变短，
故对应的布拉格级次从2增到3。

\wavequantproblem{61}{239}
一束动能 $T=180\,\mathrm{eV}$ 的
单色电子正入射镍单晶，
在与晶体表面法线成
$\theta=55^\circ$ 的方向观察到\emph{第四级}
衍射反射极大。求晶面间距 $d$。
\exsolution
电子德布罗意波长
$\lambda=h/\sqrt{2m_eT}$。
入射方向沿向内的晶面法线、反射方向相对向外法线
成 $\theta$，两束传播方向夹角是 $\pi-\theta$；
布拉格掠射角为
$\alpha=(\pi-\theta)/2$，从而
$\sin\alpha=\cos(\theta/2)$。
第四级布拉格条件
$2d\sin\alpha=4\lambda$，得
\[
 \boxed{d=\frac{2\lambda}{\cos(\theta/2)}
 =\frac{2h}{\sqrt{2m_eT}\cos(\theta/2)}
 \simeq0.21\,\mathrm{nm}}.
\]
\excheck 若误将 $55^\circ$ 直接当作对晶面的
掠射角，会产生不同的结果；此题给的是相对于\emph{法线}的出射角。

\wavequantproblem{62}{239}
动能为 $T=10\,\mathrm{keV}$ 的电子束穿过
多晶铝箔，在距箔 $\ell=10.0\,\mathrm{cm}$
的屏上产生直径 $D=3.20\,\mathrm{cm}$ 的衍射环。
若该环由\emph{三级}布拉格反射形成，
求对应晶面间距。
\exsolution
设布拉格掠射角为 $\alpha$，
衍射圆锥半角为 $2\alpha$。
屏上环半径 $D/2$，所以
\[
 \tan(2\alpha)=\frac{D}{2\ell},
 \quad\alpha=\frac12\arctan\frac{D}{2\ell}
 \simeq4.55^\circ.
\]
电子波长在 $10\,\mathrm{keV}$ 下
$\lambda\simeq h/\sqrt{2m_eT}\simeq12.3\,\mathrm{pm}$
（精确相对论值约12.2pm）。
三级布拉格关系
$2d\sin\alpha=3\lambda$，
因此
\[
 \boxed{d=\frac{3\lambda}{2\sin\alpha}
 \simeq0.23\,\mathrm{nm}}.
\]
\excheck 圆环\emph{直径}先除2求半径；
衍射方向的偏转角为 $2\alpha$，不是 $\alpha$。

\wavequantproblem{63}{239}
电子被外部电压 $U$ 加速后入射金属，
金属内势（使电子动能增加）为
$U_i=15\,\mathrm V$。
求：(a) $U=150\,\mathrm V$ 时
电子波折射率 $n$；
(b) 电子波折射率偏离1不超过
$\eta=1.0\%$ 时 $U/U_i$ 应满足的范围。
\exsolution
金属外电子动能 $T_0=eU$，
金属内的动能是 $T_i=e(U+U_i)$。
动量随 $\sqrt T$ 增大，
故德布罗意波长随 $\sqrt T$ 减小，
定义折射率
\[
 \boxed{n=\frac{\lambda_{\rm out}}{\lambda_{\rm in}}
 =\sqrt{1+\frac{U_i}{U}}}.
\]
(a) 代数值
\[
 \boxed{n=\sqrt{1.10}\simeq1.049\approx1.05}.
\]
(b) 条件 $n-1\le\eta=0.010$ 给
$1+U_i/U\le(1+\eta)^2$，故
\[
 \boxed{\frac U{U_i}\ge
 \frac1{2\eta+\eta^2}\simeq49.75\approx50}.
\]
\excheck 内势为正增加电子动能，
所以 $n>1$；当 $U\to\infty$，$n\to1$。

\wavequantproblem{64}{239}
一维无限深方势阱宽度为 $\ell$。
若仅允许整数个德布罗意\emph{半波长}
恰好装入方阱，求粒子能级。
\exsolution
边界节点条件
$\ell=n\lambda/2$，
$n=1,2,\dots$。
因此
$\lambda=2\ell/n$，动量
$p=h/\lambda=nh/(2\ell)=n\pi\hbar/\ell$。
阱内势能为0，
\[
 \boxed{E_n=\frac{p^2}{2m}
 =\frac{n^2\pi^2\hbar^2}{2m\ell^2}},
 \quad n=1,2,\dots.
\]
\excheck $n=0$ 不对应可归一化的非零定态；
基态能量严格大于0。

\wavequantproblem{65}{239}
从德布罗意波的驻波观点解释玻尔角动量量子化：
证明氢原子电子只有沿周长能容纳
\emph{整数个完整德布罗意波长}的圆轨道才满足驻波条件。
\exsolution
电子沿半径 $r$ 的圆周运动，
动量为 $p=m_ev$，
德布罗意波长为
$\lambda=h/p$。
为使波函数沿圆轨道绕行一周后具有相同相位，
要求
\[
 2\pi r=n\lambda,\qquad n=1,2,\dots.
\]
代入 $\lambda=h/p$，
\[
 2\pi rp=nh,\quad
 \boxed{L=rp=n\hbar}.
\]
这恰为玻尔的圆轨道角动量量子化条件。
\excheck 非整数圈数会导致同一点相位不连续，
无法形成稳定单值驻波；
此结论是旧量子理论的波动解释，非完整现代氢原子轨道理论。
